\section{AIME I 2002}
        \begin{enumerate}
        	\item (\textit{AIME I 2002}) Many states use a sequence of three letters followed by a sequence of three digits as their standard license-plate pattern. Given that each three-letter three-digit arrangement is equally likely, the probability that such a license plate will contain at least one palindrome (a three-letter arrangement or a three-digit arrangement that reads the same left-to-right as it does right-to-left) is $m/n$, where $m$ and $n$ are relatively prime positive integers. Find $m+n$.


 
 \textbf{Answer: }059



\textbf{Solution 1}

Consider the three-digit arrangement, $\overline{aba}$. There are $10$ choices for $a$ and $10$ choices for $b$ (since it is possible for $a=b$), and so the probability of picking the palindrome is $\frac{10 \times 10}{10^3} = \frac 1{10}$. Similarly, there is a $\frac 1{26}$ probability of picking the three-letter palindrome.  

 By the \href{https://artofproblemsolving.com/wiki/index.php?title=Principle\_of\_Inclusion-Exclusion}{Principle of Inclusion-Exclusion}, the total probability is

 $\frac{1}{26}+\frac{1}{10}-\frac{1}{260}=\frac{35}{260}=\frac{7}{52}\quad\Longrightarrow\quad7+52=\boxed{59}$

 



\textbf{Solution 2}

Using \href{https://artofproblemsolving.com/wiki/index.php?title=Complementary}{complementary} counting, we count all of the license plates that do not have the desired property.  To not be a palindrome, the first and third characters of each string must be different.  Therefore, there are $10\cdot 10\cdot 9$ three-digit non-palindromes, and there are $26\cdot 26\cdot 25$ three-letter non-palindromes.  As there are $10^3\cdot 26^3$ total three-letter three-digit arrangements, the probability that a license plate does not have the desired property is $\frac{10\cdot 10\cdot 9\cdot 26\cdot 26\cdot 25}{10^3\cdot 26^3}=\frac{45}{52}$.  We subtract this from 1 to get $1-\frac{45}{52}=\frac{7}{52}$ as our probability.  Therefore, our answer is $7+52=\boxed{59}$.

 



\textbf{Solution 3}

Note that we can pick the first and second letters/numbers freely with one choice left for the last letter/number for there to be a palindrome. Thus, the probability of no palindrome is 
 \[\frac{25}{26}\cdot \frac{9}{10}=\frac{45}{52}\] thus we have $1-\frac{45}{52}=\frac{7}{52}$ so our answer is $7+52 = \boxed{59}.$

 



\textbf{Solution 4}

We can find the probability of getting a letter and number palindrome through Solution One, which gives us $\frac{1}{26},$ and $\frac{1}{10},$ respectively. Then, we can use casework to solve the question. We begin by creating the cases:

 \begin{align*}\bullet\ \text{Case 1: The license plate includes only a letter palindrome, and no number palindrome} \\\bullet\ \text{Case 2: The license plate includes only a number palindrome, and no letter palindrome} \\\bullet\ \text{Case 3: The license plate includes both a number palindrome, and a letter palindrome}\end{align*}


 We know that the complement of these probabilities gives us the probability that the numbers and letters are NOT palindromes, so we can use that in our cases to get: 

 \begin{align}\frac{1}{26} \cdot \frac{9}{10} &= \frac{9}{260} & \text{Case 1}\\\frac{25}{26} \cdot \frac{1}{10} &= \frac{25}{260} & \text{Case 2}\\\frac{1}{26} \cdot \frac{1}{10} &= \frac{1}{260} & \text{Case 3}\end{align}

 Finally, we can add them all together to get: $\frac{9 + 25 + 1}{260} = \frac{35}{260} = \frac{7}{52} = \frac{m}{n}.$ Thus, we have $m + n = \boxed{059}.$

 


	\item (\textit{AIME I 2002}) The diagram shows twenty congruent circles arranged in three rows and enclosed in a rectangle. The circles are tangent to one another and to the sides of the rectangle as shown in the diagram. The ratio of the longer dimension of the rectangle to the shorter dimension can be written as $\frac{1}{2}\left(\sqrt{p}-q\right)$, where $p$ and $q$ are positive integers. Find $p+q$.


 
                    \begin{figure}[H]
                        \centering
                        \includegraphics[width=0.5\textwidth]{Problem/AIME_I_2002_2_Problem_Figure_1.png}
                    \end{figure}
                    


 
 \textbf{Answer: }154



\textbf{Solution}

Let the \href{https://artofproblemsolving.com/wiki/index.php?title=Radius}{radius} of the circles be $r$. The longer dimension of the rectangle can be written as $14r$, and by the \href{https://artofproblemsolving.com/wiki/index.php?title=Pythagorean\_Theorem}{Pythagorean Theorem}, we find that the shorter dimension is $2r\left(\sqrt{3}+1\right)$.

 Therefore, $\frac{14r}{2r\left(\sqrt{3}+1\right)}= \frac{7}{\sqrt{3} + 1} \cdot \left[\frac{\sqrt{3}-1}{\sqrt{3}-1}\right] = \frac{1}{2}\left(7\sqrt{3} - 7\right) = \frac{1}{2}\left(\sqrt{p}-q\right)$. Thus we have $p=147$ and $q=7$, so $p+q=\boxed{154}$.

 



\textbf{Solution 2}

Since we only care about the \textit{ratio} between the longer side and shorter side, we can set the longer side to $14$. So, this means that each of the radii is $1$. Now, we connect the radii of three circles such that they form an equilateral triangle with side length 4. Obviously, the height this triangle is $2\sqrt{3}$, and the shorter side of the triangle is therefore $2\sqrt{3}+2$ and we use simplification similar to as showed above, and we reach the result $\frac{1}{2} \cdot (\sqrt{147}-7)$ and the final answer is $147+7 = \boxed{154}$.

 


	\item (\textit{AIME I 2002}) Jane is 25 years old. Dick is older than Jane. In $n$ years, where $n$ is a positive integer, Dick's age and Jane's age will both be two-digit numbers and will have the property that Jane's age is obtained by interchanging the digits of Dick's age. Let $d$ be Dick's present age. How many ordered pairs of positive integers $(d,n)$ are possible?


 
 \textbf{Answer: }025



\textbf{Solution}

Let Jane's age $n$ years from now be $10a+b$, and let Dick's age be $10b+a$. If $10b+a>10a+b$, then $b>a$. The possible pairs of $a,b$ are:

 $(1,2), (1,3), (2,3), (1,4), (2,4), (3,4), \dots , (8,9)$

 That makes 36. But $10a+b>25$, so we subtract all the extraneous pairs: $(1,2), (1,3), (2,3), (1,4), (2,4), (1,5), (2,5), (1,6), (1,7), (1,8),$ and $(1,9)$. $36-11=\boxed{025}$

 



\textbf{Solution 2}

Start by assuming that $n < 5$ (essentially, Jane is in the 20s when their ages are 'reverses' of each other). Then we get the pairs 
 \[(61,1),(70,2),(79,3),(88,4).\] Repeating this for the 30s gives 
 \[(34,9),(43,10),(52,11),(61,12),(70,13),(79,14).\] From here, it's pretty clear that every decade we go up we get $(d,n+11)$ as a pair. Since both ages must always be two-digit numbers, we can show that each decade after the 30s, we get 1 fewer option. Therefore, our answer is $4+6+5+\dots+2+1=4+21=\boxed{025}.$~Dhillonr25

 



\textbf{Solution 3}

Note that for any "goal age", being $d + n$, there can only be $1$ ordered pair as Jane’s starting age is constant. We can work our way through Jane's possible ages when Dick can have a valid interchanged age without worrying about $n$. Referring to solution one, $b>a$.

 Starting with $a = 2$,

 


 $b = \left\{6, 7, 8, 9\right\}$ which is $4$ ages for this decade.

 
Repeating this process as $a$ gets larger, we get the sum 

 


 ${6 + 5 + 4 + 3 + 2 + 1}$.

 
Adding our numbers, we get $4 + 6 + 5 + \ldots + 2 + 1 = 4 + 21 =  \boxed{025}$. -Exit

 \textbf{Remark} The 25 ordered pairs $(d, n)$ are:(34, 9)(34, 20)(34, 31)(34, 42)(34, 53)(34, 64)(43, 10)(43, 21)(43, 32)(43, 43)(43, 54)(52, 11)(52, 22)(52, 33)(52, 44)(61, 1)(61, 12)(61, 23)(61, 34)(70, 2)(70, 13)(70, 24)(79, 3)(79, 14)(88, 4)~Puck\_0

 


	\item (\textit{AIME I 2002}) Consider the sequence defined by $a_k=\frac 1{k^2+k}$ for $k\ge 1$. Given that $a_m+a_{m+1}+\cdots+a_{n-1}=1/29$, for positive integers $m$ and $n$ with $m<n$, find $m+n$.


 
 \textbf{Answer: }840



\textbf{Solution 1}

Using \href{https://artofproblemsolving.com/wiki/index.php?title=Partial\_fraction\_decomposition}{partial fraction decomposition} yields $\dfrac{1}{k^2+k}=\dfrac{1}{k(k+1)}=\dfrac{1}{k}-\dfrac{1}{k+1}$. Thus,

 $a_m+a_{m+1}+\cdots +a_{n-1}=\dfrac{1}{m}-\dfrac{1}{m+1}+\dfrac{1}{m+1}-\dfrac{1}{m+2}+\cdots +\dfrac{1}{n-1}-\dfrac{1}{n}=\dfrac{1}{m}-\dfrac{1}{n}$

 Which means that

 $\dfrac{n-m}{mn}=\dfrac{1}{29}$

 Since we need a factor of 29 in the denominator, we let $n=29t$.* Substituting, we get

 $29t-m=mt$

 so

 $\frac{29t}{t+1} = m$

 Since $m$ is an integer, $t+1 = 29$, so $t=28$. It quickly follows that $n=29(28)$ and $m=28$, so $m+n = 30(28) = \fbox{840}$.

 \begin{verbatim}
*If $m=29t$, a similar argument to the one above implies $m=29(28)$ and $n=28$, which implies $m>n$. This is impossible since $n-m>0$.
\end{verbatim}





\textbf{Solution 2}

Note that $a_1 + a_2 + \cdots + a_i = \dfrac{i}{i+1}$. This can be proven by induction. Thus, $\sum\limits_{i=m}^{n-1} a_i = \sum\limits_{i=1}^{n-1} a_i - \sum\limits_{i=1}^{m-1} a_i = \dfrac{n-1}{n} - \dfrac{m-1}{m} = \dfrac{n-m}{mn} = 1/29$. Cross-multiplying yields $29n - 29m - mn = 0$, and adding $29^2$ to both sides gives $(29-m)(29+n) = 29^2$. Clearly, $m < n \implies 29 - m = 1$ and $29 + n = 29^2$. Hence, $m = 28$, $n = 812$, and $m+n = \fbox{840}$.

 



\textbf{Solution 3}

To solve this problem, I need to find two positive integers $m$ and $n$ where $m < n$ and the sum of sequence terms equals $\frac{1}{29}$.

 First, let me simplify $a_k = \frac{1}{k^2 + k}$ using partial fractions.$a_k = \frac{1}{k(k+1)} = \frac{1}{k} - \frac{1}{k+1}$

 Express the sum $a_m + a_{m+1} + \cdots + a_{n-1}$ using this simplification.$a_m + a_{m+1} + \cdots + a_{n-1} = \sum_{k=m}^{n-1} \left(\frac{1}{k} - \frac{1}{k+1}\right)$

 This is a telescoping series where intermediate terms cancel:$a_m + a_{m+1} + \cdots + a_{n-1} = \frac{1}{m} - \frac{1}{n}$

 Use the given condition that this sum equals $\frac{1}{29}$.$\frac{1}{m} - \frac{1}{n} = \frac{1}{29}$

 Multiplying both sides by $mn$:$n - m = \frac{mn}{29}$

 Rearranging:$29(n - m) = mn$$29n - 29m = mn$$29n - mn = 29m$$n(29 - m) = 29m$

 Solve for $n$ in terms of $m$.$n = \frac{29m}{29-m}$

 Since $n$ must be a positive integer, $29-m$ must divide $29m$ evenly.Since $29$ is prime, for $29-m$ to divide $29m$ (when $m < 29$), we need $29-m$ to divide $m$.

 This means $m = k(29-m)$ for some positive integer $k$.$m = k(29-m)$$m = 29k - km$$m(1+k) = 29k$$m = \frac{29k}{1+k}$

 For $m$ to be an integer, $1+k$ must divide $29k$.When $k = 28$, we get $m = \frac{29(28)}{29} = 28$

 Calculate $n$ using our value of $m$.$n = \frac{29(28)}{29-28} = \frac{812}{1} = 812$

 Therefore, $m + n = 28 + 812 = \boxed{840}$

 


	\item (\textit{AIME I 2002}) Let $A_1, A_2, A_3, \ldots, A_{12}$ be the vertices of a regular dodecagon. How many distinct squares in the plane of the dodecagon have at least two vertices in the set $\{A_1,A_2,A_3,\ldots,A_{12}\}$?


 
 \textbf{Answer: }183



\textbf{Solution 1}

There are 66 ways of picking two vertices. Note with any two vertices one can draw three squares (\textit{two} with the vertices forming a side, another with the vertices forming the diagonal). So so far we have $66(3)=198$ squares, but we have overcounted since some squares have their other two vertices in the dodecagon as well. All 12 combinations of two distinct vertices that form a square side only form 3 squares, and all 12 combinations of two vertices that form a square diagonal only form 6 squares. So in total, we have overcounted by $9+6=15$, and $198-15=\fbox{183}$.

 



\textbf{Solution 2}

Proceed as above to initially get 198 squares (with overcounting). Then note that any square with all four vertices on the dodecagon has to have three sides "between" each vertex, giving us a total of three squares. However, we counted these squares with all four of their sides plus both of their diagonals, meaning we counted them 6 times. Therefore, our answer is $198-3(6-1)=198-15=\boxed{183}.$

 


	\item (\textit{AIME I 2002}) The solutions to the system of equations


 
 \begin{align*} \log_{225}{x}+\log_{64}{y} = 4\\ \log_{x}{225}- \log_{y}{64} = 1 \end{align*}

are $(x_1,y_1)$ and $(x_2, y_2)$. Find $\log_{30}{(x_1y_1x_2y_2)}$.


 
 \textbf{Answer: }012



\textbf{Solution}

Let $A=\log_{225}x$ and let $B=\log_{64}y$. 

 From the first equation: $A+B=4 \Rightarrow B = 4-A$. 

 Plugging this into the second equation yields $\frac{1}{A}-\frac{1}{B}=\frac{1}{A}-\frac{1}{4-A}=1 \Rightarrow A = 3\pm\sqrt{5}$ and thus, $B=1\pm\sqrt{5}$. 

 So, $\log_{225}(x_1x_2)=\log_{225}(x_1)+\log_{225}(x_2)=(3+\sqrt{5})+(3-\sqrt{5})=6$ $\Rightarrow x_1x_2=225^6=15^{12}$. 

 And $\log_{64}(y_1y_2)=\log_{64}(y_1)+\log_{64}(y_2)=(1+\sqrt{5})+(1-\sqrt{5})=2$ $\Rightarrow y_1y_2=64^2=2^{12}$. 

 Thus, $\log_{30}\left(x_1y_1x_2y_2\right) = \log_{30}\left(15^{12}\cdot2^{12} \right) = \log_{30}\left(30^{12} \right) = \boxed{012}$. 

 
One may simplify the work by applying \href{https://artofproblemsolving.com/wiki/index.php?title=Vieta\%27s\_formulas}{Vieta's formulas} to directly find that $\log x_1 + \log x_2 = 6 \log 225, \log y_1 + \log y_2 = 2 \log 64$.

 


	\item (\textit{AIME I 2002}) 
The Binomial Expansion is valid for exponents that are not integers. That is, for all real numbers $x$, $y$, and $r$ with $|x|>|y|$, 

 $(x+y)^r=x^r+rx^{r-1}y^1+\frac{r(r-1)}2x^{r-2}y^2+\frac{r(r-1)(r-2)}{3!}x^{r-3}y^3+\cdots$

 What are the first three digits to the right of the decimal point in the decimal representation of $\left(10^{2002}+1\right)^{10/7}$?


 
 \textbf{Answer: }428



\textbf{Solution 1}

$1^n$ will always be 1, so we can ignore those terms, and using the definition ($2002 / 7 = 286$):

 
 \[(10^{2002} + 1)^{\frac {10}7} = 10^{2860}+\dfrac{10}{7}10^{858}+\dfrac{15}{49}10^{-1144}+\cdots\]

 Since the exponent of the $10$ goes down extremely fast, it suffices to consider the first few terms. Also, the $10^{2860}$ term will not affect the digits after the decimal, so we need to find the first three digits after the decimal in

 
 \[\dfrac{10}{7}10^{858}\].

 (The remainder after this term is positive by the Remainder Estimation Theorem. Since the repeating decimal of $\dfrac{10}{7}$ repeats every 6 digits, we can cut out a lot of 6's from $858$ to reduce the problem to finding the first three digits after the decimal of 

 $\dfrac{10}{7}$.

 That is the same as $1+\dfrac{3}{7}$, and the first three digits after $\dfrac{3}{7}$ are $\boxed{428}$.

 



\textbf{Solution 2}

An equivalent statement is to note that we are looking for $1000 \left\{\frac{10^{859}}{7}\right\}$, where $\{x\} = x - \lfloor x \rfloor$ is the fractional part of a number. By \href{https://artofproblemsolving.com/wiki/index.php?title=Fermat\%27s\_Little\_Theorem}{Fermat's Little Theorem}, $10^6 \equiv 1 \pmod{7}$, so $10^{859} \equiv 3^{6 \times 143 + 1} \equiv 3 \pmod{7}$; in other words, $10^{859}$ leaves a residue of $3$ after division by $7$. Then the desired answer is the first three decimal places after $\frac 37$, which are $\boxed{428}$.

 


	\item (\textit{AIME I 2002}) Find the smallest integer $k$ for which the conditions


 (1) $a_1, a_2, a_3, \ldots$ is a nondecreasing sequence of positive integers


 (2) $a_n=a_{n-1}+a_{n-2}$ for all $n>2$


 (3) $a_9=k$


 are satisfied by more than one sequence.


 
 \textbf{Answer: }748



\textbf{Solution 1}

From $(2)$, $a_9=$ $a_8+a_7=2a_7+a_6=3a_6+2a_5=5a_5+3a_4=8a_4+5a_3=13a_3+8a_2=21a_2+13a_1$ $=k$

 Suppose that $a_1=x_0$ is the smallest possible value for $a_1$ that yields a good sequence, and $a_2=y_0$ in this sequence. So, $13x_0+21y_0=k$. 

 Since $\gcd(13,21)=1$, the next smallest possible value for $a_1$ that yields a good sequence is $a_1=x_0+21$. Then, $a_2=y_0-13$. 

 By $(1)$, $a_2 \ge a_1 \Rightarrow y_0-13 \ge x_0+21 \Rightarrow y_0 \ge x_0+34 \ge 35$. So the smallest value of $k$ is attained when $(x_0,y_0)=(1,35)$ which yields $(a_1,a_2)=(1,35)$ or $(22,22)$. 

 Thus, $k=13(1)+21(35)=\boxed{748}$ is the smallest possible value of $k$.

 



\textbf{Solution 1 (Alternate)}

From $(2)$, $a_9=$ $a_8+a_7=2a_7+a_6=3a_6+2a_5=5a_5+3a_4=8a_4+5a_3=13a_3+8a_2=21a_2+13a_1$ $=k$

 Suppose that $a_2=y_0$ is the smallest possible value for $a_2$ that yields a good sequence, and $a_1=x_0$ in this sequence. So, $13x_0+21y_0=k$.

 Since $\gcd(13,21)=1$, the next smallest possible value for $a_2$ that yields a good sequence is $a_2=y_0+13$. Then, $a_1=x_0-21$. 

 $x_0-21>0$ so the smallest possible value of $x_0=22$. $y_0\ge x_0 = 22$ so the smallest value of $k$ is attained when $(x_0,y_0)=(22,22)$ which yields $(a_1,a_2)=(22,22)$ or $(22-21=1,22+13=35)$.

 Thus, $k=13(22)+21(22)=\boxed{748}$ is the smallest possible value of $k$.

 


	\item (\textit{AIME I 2002}) Harold, Tanya, and Ulysses paint a very long picket fence. Harold starts with the first picket and paints every $h$th picket; Tanya starts with the second picket and paints every $t$th picket; and Ulysses starts with the third picket and paints every $u$th picket. Call the positive integer $100h+10t+u$ $\textit{paintable}$ when the triple $(h,t,u)$ of positive integers results in every picket being painted exactly once. Find the sum of all the paintable integers.


 
 \textbf{Answer: }757



\textbf{Solution 1}

Note that it is impossible for any of $h,t,u$ to be $1$, since then each picket will have been painted one time, and then some will be painted more than once.  

 $h$ cannot be $2$, or that will result in painting the third picket twice. If $h=3$, then $t$ may not equal anything not divisible by $3$, and the same for $u$. Now for fourth and fifth pickets to be painted, $t$ and $u$ must be $3$ as well. This configuration works, so $333$ is paintable.

 If $h$ is $4$, then $t$ must be even. The same for $u$, except that it can't be $2 \mod 4$. Thus $u$ is $0 \mod 4$ and $t$ is $2 \mod 4$. Since this is all $\mod 4$, $t$ must be $2$ and $u$ must be $4$, in order for $5,6$ to be paint-able. Thus $424$ is paintable.

 $h$ cannot be greater than $4$, since if that were the case then the answer would be greater than $999$, which would be impossible for the AIME.

 Try $h=5$ for example. Then $1, 2, 3$ are covered according to the problem but $1\mod h$ covers neither $4$ nor $5$. So $t$ is forced to be $2$ for $2\mod t$ to equal $4$, and $u$ is forced to be $2$ as well for $3\mod u$ to equal $5$ (the only other way is if they both equal $1$ and that obviously double paints pickets). However, $2\mod 2$ paints all larger even numbers and $3\mod 2$ paints all larger odd numbers, which will conflict with $1 \mod 5$ many times.

 All $h \ge 5$ miss $4$ and $5$ which causes the same problem.

 Thus the sum of all paintable numbers is $\boxed{757}$.

 



\textbf{Solution 2}

Again, note that $h,t,u \neq 1$. The three conditions state that no picket number $n$ may satisfy any two of the conditions: $n \equiv 1 \pmod{h},\ n \equiv 2 \pmod{t},\ n \equiv 3 \pmod{u}$. By the \href{https://artofproblemsolving.com/wiki/index.php?title=Chinese\_Remainder\_Theorem}{Chinese Remainder Theorem}, the \href{https://artofproblemsolving.com/wiki/index.php?title=Greatest\_common\_divisor}{greatest common divisor} of any pair of the three numbers $\{h,t,u\}$ cannot be $1$ (since otherwise \href{https://artofproblemsolving.com/wiki/index.php?title=Without\_loss\_of\_generality}{without loss of generality} consider $\text{gcd}\,(h,t) = 1$; then there will be a common solution $\pmod{h \times t}$). 

 Now for $4$ to be paint-able, we require either $h = 3$ or $t=2$, but not both. 

 
                \begin{itemize}
                    \item In the former condition, since $\text{gcd}\,(h,t),\ \text{gcd}\,(h,u) \neq 1$, it follows that $3|t,u$. For $5$ and $6$ to be paint-able, we require $t = u = 3$, and it is easy to see that $333$ works.\item In the latter condition, similarly we require that $2|h,u$. All even numbers are painted. We now renumber the remaining odd pickets to become the set of all positive integers ($1,3,5, \ldots \rightarrow 1',2',3', \ldots$, where $n' = \frac{n+1}{2}$), which requires the transformation $h' = h/2,\ u' = u/2$, and with the painting starting respectively at $1',2'$. Our new number system retains the same conditions as above, except without $t$. We thus need $\text{gcd}\, (h',u') \neq 1, h',u' \neq 1$. Then for $3',4'$ to be painted, we require $h' = u' = 2$. This translates to $424$, which we see works.
                \end{itemize}


                Thus the answer is $333+424 = \boxed{757}$.

 



\textbf{Solution 3}

The three conditions state that no picket number $n$ may satisfy any two of the conditions: $n \equiv 1 \pmod{h},\ n \equiv 2 \pmod{t},\ n \equiv 3 \pmod{u}$. Note that the smallest number, $min \{ h,t,u \},$ divides the other $2$, and the next smallest divide the largest number, otherwise there is a common solution by the \href{https://artofproblemsolving.com/wiki/index.php?title=Chinese\_Remainder\_Theorem}{Chinese Remainder Theorem}. It is also a necessary condition so that it paints exactly once. Note that the smallest number can't be at least $5$, otherwise not all picket will be painted. We are left with few cases (we could also exclude $1$ as the possibility) which could be done quickly. Thus the answer is $333+424 = \boxed{757}$.

 



\textbf{Solution 4}

The wording of this problem is a bit dissatisfying and could have been improved by stating that there are at least h + t + u pickets instead of "very long". For example, a very long fence could have 5 pickets, h = 78, t = 47, u = 1. To compensate for this we should lean on the fact that the answer cannot exceed 999.

 Clearly we have to accept that there are 4 pickets or more, and the 4th picket is not reached from setting u = 1. If the 4th picket is reached from h = 3, then the 5th picket will not be reached via u = 2, or else we'd be unable to have 7 pickets. So h = 3 => t = 3 => u = 3, and 333 is a solution. We now have h > 3 for any remaining cases, and thus we only have room for 1 more case.

 Since the 4th picket must henceforth be reached via t = 2, we either have the 5th picket reached via u = 2 (requiring h to be any number exceeding the number of pickets, and we're stipulating that this is not allowed), or we have the 5th picket reached via h = 4, from which it follows that u = 4. So we have accumulated the solutions 333 and 424 and we have exhausted both the space of reasonable possibilities as well as the available space before we would exceed 999.

 


	\item (\textit{AIME I 2002}) In the diagram below, angle $ABC$ is a right angle. Point $D$ is on $\overline{BC}$, and $\overline{AD}$ bisects angle $CAB$. Points $E$ and $F$ are on $\overline{AB}$ and $\overline{AC}$, respectively, so that $AE=3$ and $AF=10$. Given that $EB=9$ and $FC=27$, find the integer closest to the area of quadrilateral $DCFG$.


 
                    \begin{figure}[H]
                        \centering
                        \includegraphics[width=0.5\textwidth]{Problem/AIME_I_2002_10_Problem_Figure_1.png}
                    \end{figure}
                    


 
 \textbf{Answer: }148



\textbf{Solution 1}

By the Pythagorean Theorem, $BC=35$. Letting $BD=x$ we can use the Angle Bisector Theorem on triangle $ABC$ to get $x/12=(35-x)/37$, and solving gives $BD=60/7$ and $DC=185/7$.

 The area of triangle $AGF$ is $10/3$ that of triangle $AEG$, since they share a common side and angle, so the area of triangle $AGF$ is $10/13$ the area of triangle $AEF$.

 Since the area of a triangle is $\frac{ab\sin{C}}2$, the area of $AEF$ is $525/37$ and the area of $AGF$ is $5250/481$.

 The area of triangle $ABD$ is $360/7$, and the area of the entire triangle $ABC$ is $210$. Subtracting the areas of $ABD$ and $AGF$ from $210$ and finding the closest integer gives $\boxed{148}$ as the answer.

 


 



\textbf{Solution 2 Bash}

By the Pythagorean Theorem, $BC=35$. From now on, every number we get will be rounded to two decimal points. Using law of cosines on AEF, EF is around 9.46. Using the angle bisector theorem, GF and DC are 7.28 and (185/7), respectively. We also know the lengths of EG and BD, so a quick Stewart's Theorem for triangles ABC and AEF gives lengths 3.76 and 14.75 for AG and AD, respectively. Now we have all segments of triangles AGF and ADC. Joy! It's time for some Heron's Formula. This gives area 10.95 for triangle AGF and 158.68 for triangle ADC. Subtracting gives us 147.73. Round to the nearest whole number and we get $\boxed{148}$.

 -jackshi2006

 



\textbf{Solution 3 Coordinate Bash}

By the Pythagorean Theorem, $BC=35$. By the Angle Bisector Theorem $BD = 60/7$ and $DC = 185/7$. We can find the coordinates of F, and use that the find the equation of line EF. Then, we can find the coordinates of G. Triangle $ADC = 1110/7$ and we can find the area triangle $AGF$ with the shoelace theorem, so subtracting that from $ADC$ gives us $\boxed{148}$ as the closest integer.

 -jackshi2006

 


 



\textbf{Solution 4 No Trig}

By the Pythagorean Theorem, $BC = 35$, and by the Angle Bisector Theorem $BD = 60/7$ and $DC = 185/7$. Draw a perpendicular from $F$ to $\overline{AE}$. Let the intersection of $F$ and $\overline{AE}$ be $H$. triangle $AHF$ is similar to $ABC$ by $AA$ similarity. thus, $AF/AC = HF/BC$. We find that $HF = 350/37$, so the area of $AEF = 525/37$. The area of triangle $AGF$ is $10/3$ that of triangle $AEG$, since they share a common side and angle, so the area of triangle $AGF$ is $10/13$ the area of triangle $AEF$, so the area of $AGF = 5250/481$. The area of triangle $ADC$ is $1110/7$, since the base is $185/7$ and the height is $12$. Thus, the area of $DCFG$ equals the area of $ADC - AGF$, or rounded to the nearest integer, $\boxed{148}$

 



\textbf{Solution 5 ONE FORMULA...THAT'S ALL! GRUBER MATH}

By the Pythagorean Theorem, $BC=35$ so $[ABC]=210$.

 Now use triangles with the same height and different bases to calculate ratios.

 Use the angle bisector theorem to calculate $DC/BD$ and thus $DC/BC=37/(37+12)$. This gives $[ADC]=\frac{37}{49}[ABC]$.

 Likewise by the angle bisector theorem for $GF/EF$, $[AGF]=\frac{10}{13}[AEF]$.

 Draw segment $FB$.

 By $AF/AC$, $[AFB]=\frac{10}{37}[ABC]$.

 By $AE/AB$, $[AEF]=\frac{1}{4}[AFB]$.

 Then $[ADC]-[AGF]=(\frac{37}{49}-\frac{10}{13}\cdot\frac{1}{4}\cdot\frac{10}{37})[ABC] \approx\boxed{148}$.

 


	\item (\textit{AIME I 2002}) Let $ABCD$ and $BCFG$ be two faces of a cube with $AB=12$. A beam of light emanates from vertex $A$ and reflects off face $BCFG$ at point $P$, which is $7$ units from $\overline{BG}$ and $5$ units from $\overline{BC}$. The beam continues to be reflected off the faces of the cube. The length of the light path from the time it leaves point $A$ until it next reaches a vertex of the cube is given by $m\sqrt{n}$, where $m$ and $n$ are integers and $n$ is not divisible by the square of any prime. Find $m+n$.


 
 \textbf{Answer: }230



\textbf{Solution 1(Easy one)}

When a light beam reflects off a surface, the path is like that of a ball bouncing. Picture that, and also imagine X, Y, and Z coordinates for the cube vertices. The coordinates will all involve 0's and 12's only, so that means that the X, Y, and Z distance traveled by the light must all be divisible by 12. Since the light's Y changes by 5 and the X changes by 7 (the Z changes by 12, don't worry about that), and 5 and 7 are relatively prime to 12, the light must make 12 reflections onto the XY plane or the face parallel to the XY plane. 

 In each reflection, the distance traveled by the light is $\sqrt{ (12^2) + (5^2) + (7^2) }$ = $\sqrt{218}$. This happens 12 times, so the total distance is $12\sqrt{218}$. $m=12$ and $n=218$, so therefore, the answer is $m+n=\boxed{230}$.

 



\textbf{Solution 2(2D and photons)}

We can use a similar trick as with reflections in 2D: Imagine that the entire space is divided into cubes identical to the one we have. Now let's follow two photons of light that start in $A$ at the same time: one of them will reflect as given in the problem statement, the second will simply fly straight through all cubes. It can easily be seen that at any moment in time the photons are in exactly the same position relative to the cubes they are inside at the moment. In other words, we can take the cube with the first photon, translate it and flip if necessary, to get the cube with the other photon.

 It follows that both photons will hit a vertex at the same time, and at this moment they will have travelled the same distance. 

 Now, the path of the second photon is simply a half-line given by the vector $(12,7,5)$. That is, the points visited by the photon are of the form $(12t,7t,5t)$ for $t\geq 0$. We are looking for the smallest $t$ such that all three coordinates are integer multiples of $12$ (which is the length of the side of the cube).

 Clearly $t$ must be an integer. As $7$ and $12$ are relatively prime, the smallest solution is $t=12$. At this moment the second photon will be at the coordinates $(12\cdot 12,7\cdot 12,5\cdot 12)$.

 Then the distance it travelled is $\sqrt{ (12\cdot 12)^2 + (7\cdot 12)^2 + (5\cdot 12)^2 } = 12\sqrt{12^2 + 7^2 + 5^2}=12\sqrt{218}$.And as the factorization of $218$ is $218=2\cdot 109$, we have $m=12$ and $n=218$, hence $m+n=\boxed{230}$.

 


	\item (\textit{AIME I 2002}) Let $F(z)=\frac{z+i}{z-i}$ for all complex numbers $z\not= i$, and let $z_n=F(z_{n-1})$ for all positive integers $n$. Given that $z_0=\frac 1{137}+i$ and $z_{2002}=a+bi$, where $a$ and $b$ are real numbers, find $a+b$.


 
 \textbf{Answer: }275



\textbf{Solution}

Iterating $F$ we get:

 
 \begin{align*} F(z) &= \frac{z+i}{z-i}\\ F(F(z)) &= \frac{\frac{z+i}{z-i}+i}{\frac{z+i}{z-i}-i} = \frac{(z+i)+i(z-i)}{(z+i)-i(z-i)}= \frac{z+i+zi+1}{z+i-zi-1}= \frac{(z+1)(i+1)}{(z-1)(1-i)}\\ &= \frac{(z+1)(i+1)^2}{(z-1)(1^2+1^2)}= \frac{(z+1)(2i)}{(z-1)(2)}= \frac{z+1}{z-1}i\\ F(F(F(z))) &= \frac{\frac{z+1}{z-1}i+i}{\frac{z+1}{z-1}i-i} = \frac{\frac{z+1}{z-1}+1}{\frac{z+1}{z-1}-1} = \frac{(z+1)+(z-1)}{(z+1)-(z-1)} = \frac{2z}{2} = z. \end{align*}


 From this, it follows that $z_{k+3} = z_k$, for all $k$. Thus $z_{2002} = z_{3\cdot 667+1} = z_1 = \frac{z_0+i}{z_0-i} = \frac{(\frac{1}{137}+i)+i}{(\frac{1}{137}+i)-i}= \frac{\frac{1}{137}+2i}{\frac{1}{137}}= 1+274i.$

 Thus $a+b = 1+274 = \boxed{275}$.

 



\textbf{Solution 2 (Very Risky, only attempt if you have very little time left)}

Solving for $z_{1}$, we have $\frac{\frac{1}{137}+2i}{\frac{1}{137}}=1+274i$. Notice that when we solve for $z_{2}$, the result is equal to some non-integral real and imaginary part, where the value of $a+b$ for $z_{2}$ is not equal to an integer. Because this is an AIME problem, the answer must be an integer, so there is most likely a pattern or cycle that $z_{n}$ goes through, and because out of the first 3 terms, only one of which has an integral value for $a+b$, you can guess that the answer is likely $1+274=\boxed{275}$

 


	\item (\textit{AIME I 2002}) In triangle $ABC$ the medians $\overline{AD}$ and $\overline{CE}$ have lengths 18 and 27, respectively, and $AB = 24$.  Extend $\overline{CE}$ to intersect the circumcircle of $ABC$ at $F$.  The area of triangle $AFB$ is $m\sqrt {n}$, where $m$ and $n$ are positive integers and $n$ is not divisible by the square of any prime.  Find $m + n$.


 
 \textbf{Answer: }063



\textbf{Solution 1}


                    \begin{figure}[H]
                        \centering
                        \includegraphics[width=0.5\textwidth]{Solution/AIME_I_2002_13_Solution_Figure_1.png}
                    \end{figure}
                    

 Applying \href{https://artofproblemsolving.com/wiki/index.php?title=Stewart\%27s\_Theorem}{Stewart's Theorem} to medians $AD, CE$, we have:

 
 \begin{align*} BC^2 + 4 \cdot 18^2 &= 2\left(24^2 + AC^2\right) \\ 24^2 + 4 \cdot 27^2 &= 2\left(AC^2 + BC^2\right)  \end{align*}


 

 Substituting the first equation into the second and simplification yields $24^2 = 2\left(3AC^2 + 2 \cdot 24^2 - 4 \cdot 18^2\right)- 4 \cdot 27^2$ $\Longrightarrow AC = \sqrt{2^5 \cdot 3 + 2 \cdot 3^5 + 2^4 \cdot 3^3 - 2^7 \cdot 3} = 3\sqrt{70}$. 

 By the \href{https://artofproblemsolving.com/wiki/index.php?title=Power\_of\_a\_Point\_Theorem}{Power of a Point Theorem} on $E$, we get $EF = \frac{12^2}{27} = \frac{16}{3}$. The \href{https://artofproblemsolving.com/wiki/index.php?title=Law\_of\_Cosines}{Law of Cosines} on $\triangle ACE$ gives 

 
 \begin{align*} \cos \angle AEC = \left(\frac{12^2 + 27^2 - 9 \cdot 70}{2 \cdot 12 \cdot 27}\right) = \frac{3}{8} \end{align*}


 

 Hence $\sin \angle AEC = \sqrt{1 - \cos^2 \angle AEC} = \frac{\sqrt{55}}{8}$ (taking the positive square root since $0^{\circ} < \angle AEC < 180^{\circ}$, so $\sin\angle AEC > 0$). Using the identity $\sin x^{\circ} \equiv \sin\left(180-x\right)^{\circ}$, we now deduce $\sin \angle AEF = \frac{\sqrt{55}}{8}$. Finally, because $\triangle AEF, BEF$ have the same height and equal bases, they have the same area, yielding 
 \[[ABF] = 2[AEF] = 2 \cdot \frac 12 \cdot AE \cdot EF \sin \angle AEF = 12 \cdot \frac{16}{3} \cdot \frac{\sqrt{55}}{8} = 8\sqrt{55},\] and so the answer is $8 + 55 = \boxed{063}$.

 



\textbf{Solution 2}

Let $AD$ and $CE$ intersect at $P$. Since medians split one another in a $2:1$ ratio, we have

 
 \begin{align*} AP = 12, PE = 9 \end{align*}


 

 This gives isosceles $APE$ and thus an easy area calculation. After extending the altitude to $PE$ and using the fact that it is also a median, we find

 
 \begin{align*} [APE] = \frac{27\sqrt{55}}{4} \end{align*}


 

 Using Power of a Point, we have

 
 \begin{align*} EF=\frac{16}{3} \end{align*}


 

 By Same Height Different Base,

 
 \begin{align*} \frac{[AFE]}{[APE]}=\frac{[AFE]}{(\frac{27\sqrt{55}}{4})}=\frac{EF}{PE}=\frac{(\frac{16}{3})}{9}=\frac{16}{27} \end{align*}


 

 Solving gives

 
 \begin{align*} [AFE] = 4\sqrt{55} \end{align*}


 

 and

 
 \begin{align*} [AFB]=2[AFE]=8\sqrt{55} \end{align*}


 

 Thus, our answer is $8+55=\boxed{063}$.

 \textbf{Short Solution: Smart Similarity} Use the same diagram as in Solution 1. Call the centroid $P$. It should be clear that $PE=9$, and likewise $AP=12$, $AE=12$. Then, $\sin \angle AEP = \frac{\sqrt{55}}{8}$. Power of a Point on $E$ gives $FE=\frac{16}{3}$, and the area of $\triangle AFB$ is $AE \cdot EF \cdot \sin \angle AEP$, which is twice the area of $\triangle AEF$ or $\triangle FEB$ (they have the same area because of equal base and height), giving $8\sqrt{55}$ for an answer of $\boxed{063}$.

 



\textbf{Solution 4 (You've Forgotten Power of a Point Exists)}

Note that, as above, it is quite easy to get that $\sin \angle AEP = \frac{\sqrt{55}}{8}$ (equate Heron's and $\frac{1}{2}ab\sin C$ to find this). Now note that $\angle FEA = \angle BEC$ because they are vertical angles, $\angle FAE = \angle ECB$, and $\angle EFA = \angle ABC$ (the latter two are derived from the inscribed angle theorem). Therefore $\Delta AEF$ ~ $\Delta CEB$ and so $FE = \frac{144}{27}$ and $\sin \angle FEA = \frac{\sqrt{55}}{8}$ so the area of $\Delta BFA$ is $8\sqrt{55}$ giving us $\boxed{063}$ as our answer. (One may just get the area via triangle similarity too--this is if you are tired by the end of test and just want to bash some stuff out--it may also serve as a useful check). 

 



\textbf{Solution 5 (Barycentric Coordinates)}

Apply barycentric coordinates on $\triangle ABC$. We know that $D=\left(0, \frac{1}{2}, \frac{1}{2}\right), E=\left(\frac{1}{2}, \frac{1}{2}, 0\right)$. We can now get the displacement vectors $\overrightarrow{AD} = \left(1, -\frac{1}{2}, -\frac{1}{2}\right)$ and $\overrightarrow{CE}=\left(-\frac{1}{2}, -\frac{1}{2}, 1\right)$. Now, applying the distance formula and simplifying gives us the two equations
 \begin{align*} 2b^2+2c^2-a^2&=1296 \\ 2a^2+2b^2-c^2&=2916. \\ \end{align*}
Substituting $c=24$ and solving with algebra now gives $a=6\sqrt{31}, b=3\sqrt{70}$. Now we can find $F$. Note that $CE$ can be parameterized as $(1:1:t)$, so plugging into the circumcircle equation and solving for $t$ gives $t=\frac{-c^2}{a^2+b^2}$ so $F=(a^2+b^2:a^2+b^2:-c^2)$. Plugging in for $a,b$ gives us $F=(1746:1746:-576)$. Thus, by the area formula, we have
 \[\frac{[AFB]}{[ABC]}= \left|\begin{matrix} 1 & 0 & 0 \\ 0 & 1 & 0 \\ \frac{97}{162} & \frac{97}{162} & -\frac{16}{81} \end{matrix}\right|=\frac{16}{81}.\]By Heron's Formula, we have $[ABC]=\frac{81\sqrt{55}}{2}$ which immediately gives $[AFB]=8\sqrt{55}$ from our ratio, extracting $\boxed{063}$.

 -Taco12

 



\textbf{Solution 6 (Law of Cosines + Stewart's theorem)}

Since $AD$ is the median, let $BD=BC=x$. Since $CE$ is a median, $AE=BE=12$. Applying Power of a Point with respect to point $E$, we see that $EF=\frac{16}{3}$. Applying Stewart's Theorem on triangles $\triangle ADC$ and $\triangle ABC$, we get that $x=3\sqrt{31}$ and $AC=3\sqrt{70}$. The area of $\triangle AFB$ is simply $\frac{1}{2} \cdot \sin{\angle FBA}\cdot FB\cdot AB$. We know $AB=24$. Also, we know that $\angle FBA = \angle FCA$. Then, applying Law of Cosines on triangle $EAC$, we get that $\cos{\angle FCA}=\frac{3\sqrt{70}}{28}$ which means that $\sin{\angle FCA=\angle FBA}=\frac{\sqrt{154}}{28}$. Then, applying Stewart's theorem on triangle $FBC$ with cevian $BE$ allows us to receive that $FB=\frac{4\sqrt{70}}{3}$. Now, plugging into our earlier area formula, we receive $\frac{1}{2} \cdot \frac{\sqrt{154}}{28} \cdot \frac{4\sqrt{70}}{3} \cdot 24 = 8\sqrt{55}.$ Therefore, the desired answer is $8+55=\boxed{063}$.

 


	\item (\textit{AIME I 2002}) A set $\mathcal{S}$ of distinct positive integers has the following property: for every integer $x$ in $\mathcal{S},$ the arithmetic mean of the set of values obtained by deleting $x$ from $\mathcal{S}$ is an integer.  Given that 1 belongs to $\mathcal{S}$ and that 2002 is the largest element of $\mathcal{S},$ what is the greatest number of elements that $\mathcal{S}$ can have?


 
 \textbf{Answer: }030



\textbf{Solution}

Let the sum of the integers in $\mathcal{S}$ be $N$, and let the size of $|\mathcal{S}|$ be $n+1$. After any element $x$ is removed, we are given that $n|N-x$, so $x\equiv N\pmod{n}$. Since $1\in\mathcal{S}$, $N\equiv1\pmod{n}$, and all elements are congruent to 1 mod $n$. Since they are positive integers, the largest element is at least $n^2+1$, the $(n+1)$th positive integer congruent to 1 mod $n$.

 We are also given that this largest member is 2002, so $2002\equiv1\pmod{n}$, and $n|2001=3\cdot23\cdot29$. Also, we have $n^2+1\le2002$, so $n \leq 44$. The largest factor of 2001 less than 45 is 29, so $n=29$ and $n+1$ $\Rightarrow{\fbox{30}}$ is the largest possible. This can be achieved with $\mathcal{S}=\{1,30,59,88,\ldots,813,2002\}$, for instance.

 


	\item (\textit{AIME I 2002}) Polyhedron $ABCDEFG$ has six faces.  Face $ABCD$ is a square with $AB = 12;$ face $ABFG$ is a trapezoid with $\overline{AB}$ parallel to $\overline{GF},$ $BF = AG = 8,$ and $GF = 6;$ and face $CDE$ has $CE = DE = 14.$  The other three faces are $ADEG, BCEF,$ and $EFG.$  The distance from $E$ to face $ABCD$ is 12.  Given that $EG^2 = p - q\sqrt {r},$ where $p, q,$ and $r$ are positive integers and $r$ is not divisible by the square of any prime, find $p + q + r.$


 
 \textbf{Answer: }163



\textbf{Solution 1}


                    \begin{figure}[H]
                        \centering
                        \includegraphics[width=0.5\textwidth]{Solution/AIME_I_2002_15_Solution_Figure_1.png}
                    \end{figure}
                    

 Let's put the polyhedron onto a coordinate plane. For simplicity, let the origin be the center of the square: $A(-6,6,0)$, $B(-6,-6,0)$, $C(6,-6,0)$ and $D(6,6,0)$. Since $ABFG$ is an isosceles trapezoid and $CDE$ is an isosceles triangle, we have symmetry about the $xz$-plane.

 Therefore, the $y$-component of $E$ is 0. We are given that the $z$ component is 12, and it lies over the square, so we must have $E(2,0,12)$ so $CE=\sqrt{4^2+6^2+12^2}=\sqrt{196}=14$ (the other solution, $E(10,0,12)$ does not lie over the square). Now let $F(a,-3,b)$ and $G(a,3,b)$, so $FG=6$ is parallel to $\overline{AB}$. We must have $BF=8$, so $(a+6)^2+b^2=8^2-3^2=55$.

 The last piece of information we have is that $ADEG$ (and its reflection, $BCEF$) are faces of the polyhedron, so they must all lie in the same plane. Since we have $A$, $D$, and $E$, we can derive this plane.* Let $H$ be the extension of the intersection of the lines containing $\overline{AG}, \overline{BF}$. It follows that the projection of $\triangle AHB$ onto the plane $x = 6$ must coincide with the $\triangle CDE'$, where $E'$ is the projection of $E$ onto the plane $x = 6$. $\triangle GHF \sim \triangle AHB$ by a ratio of $1/2$, so the distance from $H$ to the plane $x = -6$ is 
 \[\sqrt{\left(\sqrt{(2 \times 8)^2 - 6^2}\right)^2 - 12^2} = 2\sqrt{19};\] and by the similarity, the distance from $G$ to the plane $x = -6$ is $\sqrt{19}$. The altitude from $G$ to $ABCD$ has height $12/2 = 6$. By similarity, the x-coordinate of $G$ is $-6/2 = -3$. Then $G = (-6 \pm \sqrt{19}, -3, 6)$.

 Now that we have located $G$, we can calculate $EG^2$:
 \[EG^2=(8\pm\sqrt{19})^2+3^2+6^2=64\pm16\sqrt{19}+19+9+36=128\pm16\sqrt{19}.\] Taking the negative root because the answer form asks for it, we get $128-16\sqrt{19}$, and $128+16+19=\fbox{163}$.

 
                \begin{itemize}
                    \item One may also do this by vectors; $\overrightarrow{AD}\times\overrightarrow{DE}=(12,0,0)\times(-4,-6,12)=(0,-144,-72)=-72(0,2,1)$, so the plane is $2y+z=2\cdot6=12$. Since $G$ lies on this plane, we must have $2\cdot3+b=12$, so $b=6$. Therefore, $a=-6\pm\sqrt{55-6^2}=-6\pm\sqrt{19}$. So $G(-6\pm\sqrt{19},-3,6)$.
                \end{itemize}


                



\textbf{Solution 2}

We let $A$ be the origin, or $(0,0,0)$, $B = (0,0,12)$, and $D = (12,0,0)$. Draw the perpendiculars from F and G to AB, and let their intersections be X and Y, respectively. By symmetry, $FX = GY = \frac{12-6}2 = 3$, so $G = (a,b,3)$, where a and b are variables. 

 We can now calculate the coordinates of E. Drawing the perpendicular from E to CD and letting the intersection be Z, we have $CZ = DZ = 6$ and $EZ = 4\sqrt{10}$. Therefore, the x coordinate of $E$ is $12-\sqrt{(4\sqrt{10})^2-12^2}=12-\sqrt{16}=12-4=8$, so $E = (8,12,6)$.

 We also know that $A,D,E,$ and $G$ are coplanar, so they all lie on the plane $z = Ax+By+C$. Since $(0,0,0)$ is on it, then $C = 0$. Also, since $(12,0,0)$ is contained, then $A = 0$. Finally, since $(8,12,6)$ is on the plane, then $B = \frac 12$. Therefore, $b = 6$. Since $GA = 8$, then $a^2+6^2+3^2=8^2$, or $a = \pm \sqrt{19}$. Therefore, the two permissible values of $EG^2$ are $(8 \pm \sqrt{19})^2+6^2+3^2 = 128 \pm 16\sqrt{19}$. The only one that satisfies the conditions of the problem is $128 - 16\sqrt{19}$, from which the answer is $128+16+19=\boxed{163}$.

 



\textbf{Solution 3 (minimal coordinates, Apollonius)}

Denote the foot of the altitude from $E$ to $ABCD$ be $X$. Let the projection of $X$ onto $AD$ be $Y$. We seek $YD=a$. Let $E=(0, 0, 0)$. Then we get $X=(0, 0, -12)$. Because the diagram is symmetrical, $Y=(a, -6, -12)$. So, $a^2+6^2+12^2=14^2 \rightarrow a=4$. We find $EA=2\sqrt{61}$.

 Extend $EG$ and $FE$ to meet the plane $z=0$. Since $EGAD$ is a quadrilateral and all on a plane, then the extension of $EG$ and $FE$ will meet the lines $AD$ and $BC$, respectively. Call these intersections $A'$ and $B'$. Let $EA'=a, AA'=b$.

 Using the Law of Cosines on $\triangle EAD$ gives $\cos(\angle EAD)=\frac{4}{\sqrt{61}}$.Using Law of Cosines on $\triangle EA'D$ gives the equation $a^2=b^2+244+16b$. Now, using Apollonius' Theorem on the same triangle gives $a^2=2b^2+232$.Equating the two gives $b^2-16b-12=0$. Solving gives us $b=8-2\sqrt{19}, b^2=140-32\sqrt{19}$.

 Finally, plugging into either expression for $a$ gives $a^2=512-64\sqrt{19}$.Since $FG=\frac{1}{2}A'B'$ and is parallel to $A'B'$, by the midpoint theorem, 
 \[EG=\frac{1}{2} A'B' \rightarrow EG^2=\frac{1}{4}A'B'^2 \rightarrow = 128-16\sqrt{19}\].

 Then $128+16+19=163$.

 -RyanZhu817

 


\end{enumerate}